\documentclass[11pt]{article}

\usepackage[margin=1in]{geometry}
\usepackage{hyperref}
\usepackage{amsmath}
\usepackage{enumitem}
\usepackage{listings}
\usepackage{xcolor}

\hypersetup{
    colorlinks=true,
    linkcolor=blue,
    urlcolor=blue,
    citecolor=blue
}

\title{CS 300 -- Week 4 Resources\\\large Hash Tables: Introduction and Collision Handling}
\author{Jose de Lima}
\date{\today}

\begin{document}
\maketitle

\section*{Overview}
Week 4 introduces \textbf{hash tables}, the data structure that gives us
\emph{average-case} $O(1)$ insert, search, and delete by translating keys
into array indices through a \textbf{hash function}. The resources below
cover two things: the core idea of hashing, and the most common strategy
for coping with the unavoidable reality of \textbf{collisions} --
\textbf{separate chaining}.

\section{Hashing -- Introduction}
\textbf{Resource:} \href{https://www.youtube.com/watch?v=wWgIAphfn2U}{Hashing Set 1 -- Introduction (GeeksforGeeks, YouTube)}

\subsection*{The Problem Hashing Solves}
Arrays give us $O(1)$ access \emph{by index}, but most real-world lookups
are by \emph{key} (a student ID, a bid number, a string). Walking an array
to find a matching key is $O(n)$; a balanced BST trims that to
$O(\log n)$. Hashing aims for $O(1)$ on average by converting the key
itself into an array index.

\subsection*{The Three Ingredients}
\begin{enumerate}[leftmargin=*]
    \item \textbf{Keys} -- the identifiers we want to look items up by.
    \item \textbf{A hash function} $h(key)$ -- deterministic, fast, and
          distributes keys as uniformly as possible across the table.
    \item \textbf{A bucket array} of size $m$ -- the storage; the index into
          it is $h(key) \bmod m$.
    \end{enumerate}

\subsection*{Properties of a Good Hash Function}
\begin{itemize}[leftmargin=*]
    \item \textbf{Deterministic:} the same key always maps to the same slot.
    \item \textbf{Uniform distribution:} minimizes clustering, which in turn
          minimizes collisions.
    \item \textbf{Fast to compute:} ideally $O(1)$ in the size of the key.
    \item \textbf{Uses every part of the key:} so small changes to the key
          scatter across the table rather than colliding.
\end{itemize}

\subsection*{Collisions Are Unavoidable}
By the pigeonhole principle, if the universe of possible keys is larger
than the table size $m$ (which it essentially always is), two distinct
keys will eventually hash to the same slot. A hash-table implementation
must define a \textbf{collision-resolution strategy}. The two broad
families are:
\begin{itemize}[leftmargin=*]
    \item \textbf{Separate chaining} -- each slot holds a secondary
          structure (usually a linked list) of all keys that hashed there.
    \item \textbf{Open addressing} -- on a collision, probe other slots in
          the same array (linear probing, quadratic probing, double
          hashing).
\end{itemize}

\subsection*{Load Factor}
The \textbf{load factor} $\alpha = n / m$ (items over buckets) is the
single most important performance knob. As $\alpha$ rises, collisions
rise and lookups slow down. Most implementations \textbf{resize} (grow
the bucket array and \emph{rehash} everything into it) when $\alpha$
crosses a threshold -- commonly $0.75$.

\section{Hash Tables with Chaining}
\textbf{Resource:} \href{https://www.youtube.com/watch?v=9IXCdJA8O3c}{CS 260 Lesson 5 -- Hash Tables with Chaining (YouTube)}

\subsection*{The Idea}
In separate chaining, each bucket is itself a container -- classically a
singly linked list -- that stores \emph{every} key-value pair whose hash
lands in that slot. Collisions don't overwrite anything; they just
lengthen a chain.

\begin{center}
\texttt{[0] -> (key\_a, val) -> (key\_g, val) -> NULL}\\
\texttt{[1] -> NULL}\\
\texttt{[2] -> (key\_b, val) -> NULL}\\
\texttt{[3] -> (key\_c, val) -> (key\_f, val) -> (key\_k, val) -> NULL}
\end{center}

\subsection*{Core Operations}
Let $h(key)$ be the hash function and let \texttt{table[i]} be the
linked list in bucket $i$.

\paragraph{Insert(key, value).}
\begin{enumerate}[leftmargin=*]
    \item Compute $i = h(key) \bmod m$.
    \item Walk \texttt{table[i]}. If a node with this key exists, update
          its value (upsert).
    \item Otherwise append (or prepend) a new node at the front of
          \texttt{table[i]}.
\end{enumerate}

\paragraph{Search(key).}
\begin{enumerate}[leftmargin=*]
    \item Compute $i = h(key) \bmod m$.
    \item Walk \texttt{table[i]} until a node with a matching key is
          found -- return its value -- or the end is reached -- return
          ``not found.''
\end{enumerate}

\paragraph{Remove(key).}
\begin{enumerate}[leftmargin=*]
    \item Compute $i = h(key) \bmod m$.
    \item Walk \texttt{table[i]}, keeping track of the predecessor.
    \item When the matching node is found, unlink it from the chain and
          free it.
\end{enumerate}

\subsection*{Pseudocode}
\begin{lstlisting}[basicstyle=\ttfamily\small,frame=single]
function insert(key, value):
    i = hash(key) % m
    for node in table[i]:
        if node.key == key:
            node.value = value        # upsert
            return
    prepend (key, value) to table[i]

function search(key):
    i = hash(key) % m
    for node in table[i]:
        if node.key == key:
            return node.value
    return NOT_FOUND

function remove(key):
    i = hash(key) % m
    walk table[i] tracking (prev, cur)
    when cur.key == key:
        unlink cur from the chain
        delete cur
\end{lstlisting}

\subsection*{Complexity}
Let $n$ be the number of items, $m$ the number of buckets,
$\alpha = n/m$.
\begin{itemize}[leftmargin=*]
    \item \textbf{Average case} (good hash, reasonable $\alpha$):
          $O(1 + \alpha)$ for insert, search, remove -- effectively
          $O(1)$ when the table is kept at a bounded load factor.
    \item \textbf{Worst case}: $O(n)$ -- every key collides into the
          same chain.
    \item \textbf{Space}: $O(n + m)$.
\end{itemize}

\subsection*{Separate Chaining vs. Open Addressing}
\begin{center}
\begin{tabular}{l|l|l}
\textbf{Property} & \textbf{Chaining} & \textbf{Open Addressing} \\ \hline
Storage location    & External linked list per bucket & All entries inside the table  \\
Handles $\alpha > 1$? & Yes (chains just grow)        & No -- table must be resized    \\
Cache locality      & Worse (pointer chasing)         & Better (contiguous array)     \\
Deletion            & Trivial (unlink a node)         & Requires tombstones           \\
Memory overhead     & Extra pointers per node         & None per entry                \\
Worst-case lookup   & $O(n)$                          & $O(n)$                        \\
\end{tabular}
\end{center}

\subsection*{Practical Notes}
\begin{itemize}[leftmargin=*]
    \item A prime-sized table tends to scatter keys better under
          \texttt{mod}, reducing clustering when the hash function is
          weak.
    \item Track the load factor on every insert; rehash when it crosses
          the threshold (e.g., $0.75$) by allocating a larger table and
          re-inserting every entry.
    \item If chains grow long under adversarial input, some libraries
          swap the linked list for a balanced tree once a chain exceeds
          a threshold (Java's \texttt{HashMap} does this).
\end{itemize}

\section*{Key Takeaways}
\begin{itemize}[leftmargin=*]
    \item A hash table is an array plus a hash function plus a collision
          strategy -- all three have to be right for the $O(1)$ promise
          to hold.
    \item Separate chaining is the simplest collision strategy: each
          bucket is a linked list, so collisions just lengthen a chain
          instead of cascading through the table.
    \item Performance tracks the load factor $\alpha$. Keep it bounded
          (resize when needed) and expected operations stay $O(1)$.
\end{itemize}

\section*{Links Summary}
\begin{itemize}[leftmargin=*]
    \item Hashing -- Introduction: \url{https://www.youtube.com/watch?v=wWgIAphfn2U}
    \item Hash Tables with Chaining: \url{https://youtu.be/9IXCdJA8O3c}
\end{itemize}

\end{document}
